% ==========================================================================
%  CV — Clément Noël
% --------------------------------------------------------------------------
%  Compilation : pdflatex cv_clement_noel.tex   (ou Overleaf, compilateur pdfLaTeX)
%  Paquets utilisés : tous inclus dans TeX Live, MiKTeX et Overleaf.
%
%  ORGANISATION
%   1. Réglages (couleurs, coordonnées)   ← à modifier en premier
%   2. Mise en page et commandes du modèle
%   3. Contenu du CV
%
%  COMMANDES DU MODÈLE
%   \rubrique{Titre}                          titre de section
%   \experience{période}{poste}{entreprise}{détails}
%   \formation{période}{diplôme}{établissement}{détails}
%   \competence{nom}{niveau de 0 à 5}
% ==========================================================================

\documentclass[10pt,a4paper]{article}

\usepackage[utf8]{inputenc}
\usepackage[T1]{fontenc}
\usepackage[french]{babel}
\usepackage[default]{sourcesanspro}   % police sans empattement
\usepackage[a4paper,margin=1.3cm,top=1.1cm,bottom=1.1cm]{geometry}
\usepackage{xcolor}
\usepackage{graphicx}
\usepackage{tabularx}
\usepackage{enumitem}
\usepackage{fontawesome5}
\usepackage[hidelinks,pdfusetitle]{hyperref}


% ==========================================================================
%  1. RÉGLAGES
% ==========================================================================
\definecolor{accent}{HTML}{1F3A5F}     % bleu nuit : titres et jauges
\definecolor{texte}{HTML}{222222}
\definecolor{discret}{HTML}{6B6B6B}    % dates, lieux
\definecolor{clair}{HTML}{D5DAE1}      % segments vides des jauges

\newcommand{\nom}{Clément Noël}
\newcommand{\poste}{Automaticien — Automatisme \& systèmes embarqués}
\newcommand{\ville}{33160 Pessac}
\newcommand{\email}{clementnoelrenard2001@ik.me}
\newcommand{\linkedin}{clement-noel-renard}
\newcommand{\siteweb}{clementnoelrenard.github.io}
\newcommand{\photo}{../images/MOI3.jpg}   % laisser tel quel si pas de photo

\hypersetup{
    pdftitle    = {Curriculum vitæ — Clément Noël},
    pdfauthor   = {Clément Noël},
    pdfsubject  = {Curriculum vitæ},
    pdfkeywords = {automatisme, informatique industrielle, BUT GEII, BTS CRSA}
}


% ==========================================================================
%  2. MISE EN PAGE ET COMMANDES DU MODÈLE
% ==========================================================================
\pagestyle{empty}
\setlength{\parindent}{0pt}
\setlength{\tabcolsep}{0pt}
\color{texte}
\setlist[itemize]{leftmargin=1em, label={\color{accent}\textbullet}, nosep, topsep=1pt}

% Titre de section : petites capitales espacées, filet en dessous
\newcommand{\rubrique}[1]{%
    \vspace{7pt}%
    {\color{accent}\large\bfseries\MakeUppercase{#1}}\par
    \vspace{-3pt}{\color{accent}\rule{\linewidth}{0.6pt}}\par
    \vspace{3pt}%
}

% Bloc générique : période à droite, titre en gras, organisme en italique
\newcommand{\bloc}[4]{%
    \begin{tabularx}{\linewidth}{@{}X r@{}}
        \textbf{#2} & {\color{discret}\small #1} \\
    \end{tabularx}\par
    {\color{accent}\itshape #3}\par
    \ifx\relax#4\relax\else{\small #4}\fi
    \vspace{5pt}%
}
\newcommand{\experience}[4]{\bloc{#1}{#2}{#3}{#4}}
\newcommand{\formation}[4]{\bloc{#1}{#2}{#3}{#4}}

% Jauge de compétence sur 5 segments
\newcommand{\plein}{{\color{accent}\rule{8pt}{3.2pt}}\hspace{1.5pt}}
\newcommand{\vide}{{\color{clair}\rule{8pt}{3.2pt}}\hspace{1.5pt}}
\newcommand{\jauge}[1]{%
    \ifcase#1 \vide\vide\vide\vide\vide
    \or \plein\vide\vide\vide\vide
    \or \plein\plein\vide\vide\vide
    \or \plein\plein\plein\vide\vide
    \or \plein\plein\plein\plein\vide
    \else \plein\plein\plein\plein\plein
    \fi
}
\newcommand{\competence}[2]{\small #1 \hfill \jauge{#2}\par\vspace{1.5pt}}
\newcommand{\domaine}[1]{\vspace{3pt}{\bfseries\small #1}\par\vspace{1pt}}

% Élément de coordonnées : icône + texte
\newcommand{\coord}[2]{{\color{accent}#1}\,\,#2}


% ==========================================================================
%  3. CONTENU DU CV
% ==========================================================================
\begin{document}

% ---------------------------------------------------------------- En-tête --
\begin{minipage}[c]{0.80\linewidth}
    {\fontsize{26}{28}\selectfont\bfseries\color{accent}\nom}\par
    \vspace{4pt}
    {\Large \poste}\par
    \vspace{8pt}
    \small
    \coord{\faMapMarker*}{\ville}\quad
    \coord{\faEnvelope}{\href{mailto:\email}{\email}}\quad
    \coord{\faLinkedin}{\href{https://www.linkedin.com/in/\linkedin/}{\linkedin}}\quad
    \coord{\faGlobe}{\href{https://\siteweb/}{\siteweb}}
\end{minipage}%
\hfill
\begin{minipage}[c]{0.17\linewidth}
    \raggedleft
    \IfFileExists{\photo}{\includegraphics[width=\linewidth]{\photo}}{}
\end{minipage}

\vspace{10pt}
{\small Double diplômé BTS CRSA et BUT GEII (parcours Automatisme et Informatique Industrielle),
je conçois, programme et mets en service des systèmes automatisés, de l'armoire électrique à la
supervision. Passionné par la protection de la vie privée et l'autonomie numérique, je contribue
à Wikipédia, OpenStreetMap et à la documentation francophone de Debian.\par}
\vspace{4pt}
{\small\bfseries\color{accent} En recherche d'une alternance 2026 – 2029 pour le cycle ingénieur
industriel de l'ESTIA (Bidart).\par}

% ------------------------------------------------------- Colonne gauche --
\vspace{2pt}
\begin{minipage}[t]{0.62\linewidth}

\rubrique{Expérience professionnelle}

\experience{mars 2025 – juil. 2026}{Alternant automaticien}{INEO Aquitaine (Equans) — conception industrielle}{%
\begin{itemize}
    \item Co-conception du rétrofit d'interfaces opérateur Schneider Electric
    \item Installation d'une baie industrielle avec KVM et automate Siemens
    \item Rigueur, gestion de la charge de travail et travail en équipe
\end{itemize}}

\experience{avr. 2024 – juil. 2024}{Stagiaire automaticien — bureau d'études}{ISP Aquitaine — conception industrielle}{%
\begin{itemize}
    \item Conception et mise en œuvre d'une supervision OPC UA
    \item Intégration d'une cellule robotique ABB
\end{itemize}}

\experience{avr. 2021 – juil. 2021}{Stagiaire électricien de maintenance}{Lacosse Emballage — scierie}{%
\begin{itemize}
    \item Restauration complète et mise aux normes d'une machine de découpe de bois
    \item Projet complet : électricité, mécanique, pneumatique, programmation
\end{itemize}}

\experience{janv. 2020 – févr. 2020}{Stagiaire vendeur conseil}{Yess Électrique — matériel électrique pour professionnels}{%
Conseil technique à partir d'un cahier des charges, relation client.}

\rubrique{Formation}

\formation{2023 – 2026}{BUT GEII — Automatisme et Informatique Industrielle}{IUT de Bordeaux, Gradignan}{%
Programmation, électronique, informatique, énergie, automatisme. Trois projets en conditions d'entreprise.}

\formation{2020 – 2023}{BTS CRSA — Conception et Réalisation de Systèmes Automatiques}{Lycée Val de Garonne, Marmande}{%
Mécanique, informatique, automatisme. Projet d'un an : conception complète d'un process et programmation robotique (Roboguide).}

\rubrique{Projets}

\begin{itemize}
    \item \textbf{SAE S4 MugOchaud} — maintien en température d'un mug : armoire, automate (LADDER), maintenance
    \item \textbf{SAE S3 Robot Mini-Sumo} — robot autonome, détection d'adversaire par capteurs VL53L0X
    \item \textbf{Projet tuteuré Cave à vin} — régulation PID température/humidité sous TIA Portal
    \item \textbf{BTS CRSA} — pupitre opérateur d'une machine de tri de ronds de serviette
\end{itemize}

\end{minipage}%
\hfill
% ------------------------------------------------------- Colonne droite --
\begin{minipage}[t]{0.33\linewidth}

\rubrique{Compétences}

\domaine{Programmation industrielle}
\competence{EcoStruxure / Vijeo Designer}{5}
\competence{Unity Pro XL}{5}
\competence{TIA Portal V21}{4}

\domaine{Électricité \& électronique}
\competence{Mise aux normes, entretien}{5}
\competence{Schémas XRelais, SEE Electrical}{4}
\competence{Dépannage de cartes}{4}
\competence{Schémas EasyEDA, KiCad}{3}

\domaine{Conception \& fabrication}
\competence{Impression 3D, découpe laser}{5}
\competence{Inventor, Fusion 360}{4}
\competence{Statique, cinématique}{3}
\competence{SolidWorks}{2}

\domaine{Informatique}
\competence{Windows, Linux, VMware}{5}
\competence{LLM, protection des données}{5}
\competence{Docker, VLAN, VPN}{3}
\competence{C++, Python, HTML}{2}
\competence{MATLAB, GNU Octave}{1}

\rubrique{Diplômes}

{\small
\begin{tabularx}{\linewidth}{@{}X r@{}}
    BUT GEII                 & 2026 \\
    DUT GEII                 & 2025 \\
    BTS CRSA                 & 2023 \\
    Bac Pro MELEC            & 2020 \\
    BEP Électrotechnique     & 2019 \\
    Brevet d'Initiation Aéronautique & 2019 \\
\end{tabularx}\par}

\rubrique{Centres d'intérêt}

{\small
Électronique, aéronautique et spatial (Kerbal Space Program), automatisation sur Factorio,
serveur auto-hébergé, conception 3D, trail dans les Pyrénées. Participation à une olympiade
de robotique.\par}

\end{minipage}

\end{document}
